%=============================================================================
% WAVE 3 ITERATION 5: Patent 8 (Communications) + Patent 9 (Navigation)
%
% DFD Research Programme -- March 2026
% Iteration 5 of 20 -- M_eff correction propagation, honest DSN comparison,
%   secrecy survival, sub-mm navigation audit, n_psi=sqrt(2) systematic
%
% W3i5 MANDATE (P8):
%   1. Recalculate Earth-Mars data rate with corrected M_eff = 3.01e6 kg
%   2. Does inherent secrecy survive the M_eff correction?
%   3. Corrected channel capacity
%   4. Honest comparison with EM deep-space comms (DSN)
%
% W3i5 MANDATE (P9):
%   5. M_eff impact on sub-mm navigation claims
%   6. Corrected geological detection depth/resolution
%   7. Does 8.7 nm underwater positioning survive?
%   8. Impact of n_psi=sqrt(2) systematic on MOND-zone positioning
%=============================================================================
\documentclass[aps,prd,twocolumn,superscriptaddress,nofootinbib]{revtex4-2}

\usepackage{amsmath,amssymb,amsthm}
\usepackage{mathtools}
\usepackage{physics}
\usepackage{siunitx}
\usepackage{booktabs}
\usepackage{hyperref}
\usepackage{xcolor}
\usepackage{array}
\usepackage{multirow}
\usepackage{empheq}

\newcommand{\DFD}{\textsc{dfd}}
\newcommand{\psifield}{\psi}
\newcommand{\astar}{a_*}
\newcommand{\azero}{a_0}
\newcommand{\mufunction}{\mu}
\newcommand{\order}[1]{\mathcal{O}(#1)}
\newcommand{\SNR}{\mathrm{SNR}}
\newcommand{\BER}{\mathrm{BER}}
\newcommand{\Cs}{C_{\rm s}}
\newcommand{\Cw}{C_{\rm w}}
\newcommand{\Ie}{I_{\rm Eve}}
\newcommand{\PSF}{\mathrm{PSF}}
\newcommand{\CRLB}{\mathrm{CRLB}}
\newcommand{\drho}{\Delta\rho}
\newcommand{\km}{\,\mathrm{km}}
\DeclareMathOperator{\Tr}{Tr}

% Colour-coded verdict boxes
\newcommand{\confirmed}[1]{\textcolor{blue}{\textbf{CONFIRMED:} #1}}
\newcommand{\corrected}[1]{\textcolor{red}{\textbf{CORRECTED:} #1}}
\newcommand{\caution}[1]{\textcolor{orange}{\textbf{CAUTION:} #1}}
\newcommand{\honest}[1]{\textcolor{violet}{\textbf{HONEST ASSESSMENT:} #1}}
\newcommand{\killed}[1]{\textcolor{red!70!black}{\textbf{KILLED:} #1}}

\newtheorem{theorem}{Theorem}
\newtheorem{proposition}[theorem]{Proposition}
\newtheorem{lemma}[theorem]{Lemma}
\newtheorem{corollary}[theorem]{Corollary}
\newtheorem{finding}{Finding}
\newtheorem{correction}{Correction}
\newtheorem{definition}{Definition}
\newtheorem{warning}{Warning}
\newtheorem{remark}{Remark}

\begin{document}

\title{Wave 3 Iteration 5: Patent 8 (Communications) + Patent 9 (Navigation)\\
\textit{M$_{\rm eff}$ Correction Propagation $\cdot$ Honest DSN Comparison\\
Secrecy Survival $\cdot$ Sub-mm Navigation Audit\\
n$_\psi = \sqrt{2}$ Systematic on MOND-Zone Positioning}}

\author{DFD Research Programme --- Agent 8/9, Iteration 5}
\date{20 March 2026}

\begin{abstract}
This Iteration~5 combined update for Patents~8 and~9 propagates the
critical W3i4 correction $M_{\rm eff} = 3.01\times10^6$~kg (replacing
the placeholder $M_\psi = 10^{-6}$~kg) through both patent chains,
and performs honest audits mandated by the W3i4 findings.

\textbf{Patent 8 (Communications):}
(i)~The W3i4 Earth-Mars link budget already used $M_s = 8.3\times10^{23}$~kg
(planetary-mass source) because the raw SNR with $M_s = 10^6$~kg was
$\SNR \approx 2.4\times10^{-47}$ (zero rate). The $M_{\rm eff}$ correction
does not alter the link budget: the source mass in the link budget is the
\emph{transmitter} mass, not the DFD mode mass. The Earth-Mars opposition rate
remains $\sim 12.2$~kbits/s ($\eta = 1$) or $14.0$~kbits/s (MRC, $\eta = 3$),
but only with a planetary-mass transmitter.
(ii)~Inherent secrecy \textbf{survives}: the secrecy capacity
$C_s = B\log_2((1+\SNR_B)/(1+\SNR_E)) > 0$ depends only on the
\emph{geometric} channel structure (Bob inside corridor, Eve outside),
not on the mode mass.
(iii)~Honest DSN comparison: NASA DSN achieves 0.5--6~Mbits/s at Mars
opposition using 34--70~m antennas at X/Ka band. The DFD psi-channel at
$\sim 14$~kbits/s is $\sim 35$--$430\times$ slower, and requires a
planetary-mass transmitter. The psi-channel advantage is \emph{not} data rate
but (a)~zero conjunction blackout, (b)~inherent secrecy, and
(c)~$1/d^2$ not $1/d^4$ power scaling.

\textbf{Patent 9 (Navigation):}
(iv)~The $M_{\rm eff}$ correction does not affect P9 positioning because
positioning uses \emph{ambient} $\psi$-field gradients from natural
masses (Earth, geological bodies), not from a DFD-generated mode.
All sub-mm and 8.7~nm claims are \textbf{unchanged}.
(v)~Geological detection: the minimum detectable anomaly formula
$M_{\rm anom}^{\min} = 5\sigma_\psi c^2 d / G$ is independent of
$M_{\rm eff}$ (it depends on sensor noise floor, not mode mass).
All W3i4 geological results \textbf{survive}.
(vi)~The $n_\psi = \sqrt{2}$ correction affects MOND-zone positioning
at altitudes $\gtrsim 3200$~km where $|\nabla\psi| \lesssim a_0$.
The psi-field propagation speed $c_s = c/n_\psi = c/\sqrt{2}$ introduces
a systematic $\sqrt{2}$ correction to ranging-based position estimates in
the MOND zone. For surface and near-surface operations (the primary P9 use
case), $n_\psi \approx 1$ and the correction is negligible
($\lesssim 10^{-10}$).
\end{abstract}

\maketitle
\tableofcontents

%=============================================================================
\section{Overview: What the $M_{\rm eff}$ Correction Means for P8 and P9}
\label{sec:overview}
%=============================================================================

\subsection{The W3i4 $M_{\rm eff}$ Correction (Recalled)}

W3i4 P11 established that the DFD gravitational mode mass is:
\begin{equation}
M_\psi^{\rm grav} = \frac{\mu_0 V_{\rm cav}}{4\pi G}
\approx 3.01\times10^6\,\mathrm{kg}
\label{eq:Meff}
\end{equation}
for a SQUID-scale cavity. The previously used placeholder
$M_\psi = 10^{-6}$~kg was never derived from DFD first principles.
The ratio $M_\psi^{\rm grav}/M_\psi^{\rm placeholder} = 3\times10^{12}$,
making all \emph{absolute sensitivity estimates} that used the placeholder
$\sim 10^6\times$ over-optimistic (the square-root of the mass ratio
enters linearly in signal amplitudes: $\delta\psi \propto M_s/d$,
so if $M_s$ was underestimated by $10^{12}$, signals were overestimated
by $10^6$ in the coupling regime).

\subsection{Key Question: Does $M_{\rm eff}$ Enter P8 or P9 Physics?}

\begin{finding}[$M_{\rm eff}$ Relevance Triage]
\label{find:triage}
The gravitational mode mass $M_\psi^{\rm grav}$ governs:
\begin{enumerate}
\item The response of a \emph{cavity-confined} psi-mode to EM driving
  (P2, P11, P15 --- active generation/detection).
\item The noise floor of a SQUID-based psi-detector coupling to a
  cavity mode (P11).
\end{enumerate}

It does \textbf{not} govern:
\begin{enumerate}
\item[(a)] The free-space propagation of the psi-field from a
  macroscopic mass (P8 link budget uses $M_s = $ transmitter mass,
  not cavity mode mass).
\item[(b)] The ambient psi-field gradient from Earth/geological bodies
  used for P9 positioning (these are Newtonian: $\psi = GM/c^2 r$).
\item[(c)] The sensor noise floor of atom interferometers (P6/P9),
  which is set by shot noise and quantum projection noise, not by
  the DFD mode mass.
\item[(d)] The secrecy capacity, which is a geometric/information-theoretic
  property of the channel.
\end{enumerate}
\end{finding}

\begin{correction}[Scope of $M_{\rm eff}$ Impact on P8/P9]
\label{corr:scope}
\textbf{P8}: The Earth-Mars link budget (W3i4 Sec.~VI) already identified
that $M_s = 10^6$~kg gives $\SNR = 2.4\times10^{-47}$ (zero rate), and
that a planetary-mass source ($M_s \sim 8.3\times10^{23}$~kg) is required.
The $M_{\rm eff}$ correction does not change this conclusion --- it
reinforces it. The link budget equation
\begin{equation}
C = B\log_2\!\left(1 + \eta \frac{G^2 M_s^2 \xi_0^2}{c^4 d^2 N_0 B}\right)
\label{eq:link-budget-recalled}
\end{equation}
uses the \emph{transmitter's total mass} $M_s$, not the DFD cavity mode
mass. The mode mass enters only if one attempts to \emph{generate} $\psi$
modulation from a cavity (P15 domain), not in the propagation or reception.

\textbf{P9}: All positioning and geological detection formulas use
$\psi_0 = GM_\oplus/(c^2 R_\oplus)$ and
$\delta\psi = G M_{\rm anom}/(c^2 d)$, which are determined by the
\emph{physical masses} of Earth and geological bodies, not by the
DFD mode mass. The sensor noise $\sigma_\psi$ is determined by P6
quantum sensor performance. Neither quantity involves $M_{\rm eff}$.

\corrected{P8 and P9 are \textbf{unaffected} by the $M_{\rm eff}$
correction. The correction impacts P2, P11, and P15 (active
generation/detection), not P8 (free-space communication) or P9
(passive navigation/geological sensing).}
\end{correction}

%=============================================================================
\section{P8 Task 1: Earth-Mars Data Rate --- Recalculated}
\label{sec:earth-mars}
%=============================================================================

\subsection{Link Budget with Physical Source Mass}

The W3i4 link budget (Finding~8 of W3i4 P8) established:
\begin{equation}
\boxed{
C_{\rm E-M}(d) = B\log_2\!\left(1 + \eta \cdot
\frac{G^2 M_s^2 \xi_0^2}{c^4 d^2 N_0 B}\right),
}
\label{eq:link-budget}
\end{equation}
with parameters: $B = 10^3$~Hz, $\xi_0 = 10^{-3}$,
$N_0 = 2.5\times10^{-38}$~m$^2$s$^{-2}$Hz$^{-1}$, $\eta = 3$ (SIMO MRC).

The required source mass for $\SNR_{\rm opp} = 4807$ at
$d_{\rm opp} = 5.6\times10^{10}$~m was derived as
$M_s \approx 8.3\times10^{23}$~kg (Venus-mass).

\subsection{Does $M_{\rm eff} = 3.01\times10^6$~kg Change Anything?}

No. The link budget source mass $M_s$ is the total mass of the
transmitting body whose mass-energy is modulated. For deep-space
communication, this would be:
\begin{itemize}
\item A natural body (asteroid, dwarf planet) with mass modulation
  via surface mass redistribution.
\item A very large artificial structure (space station) with modulated
  mass distribution.
\end{itemize}

The DFD mode mass $M_{\rm eff}$ governs the cavity response in a
P15-type generator, not the free-space field from a macroscopic source.
A P15 cavity does not \emph{amplify} the gravitational field; it
couples EM energy into the psi-mode at rate set by $M_{\rm eff}$.

\begin{warning}[P15 as Transmitter]
\label{warn:p15-tx}
If one attempted to use a P15 psi-field generator as the transmitter
(rather than modulating a macroscopic mass), the radiated psi-power
would be suppressed by $M_{\rm eff}$:
\begin{equation}
P_\psi^{\rm P15} \sim \frac{(\lambda-1)^2 P_{\rm EM}}{M_{\rm eff} c^2}
\ll P_{\rm EM}.
\label{eq:P15-power}
\end{equation}
With $M_{\rm eff} = 3.01\times10^6$~kg and current bounds
$|\lambda-1| < 4.9\times10^{-7}$, the radiated psi-power from a P15
generator at $P_{\rm EM} = 1$~MW is:
\begin{equation}
P_\psi^{\rm P15} \sim \frac{(4.9\times10^{-7})^2 \times 10^6}
{3.01\times10^6 \times 9\times10^{16}}
\approx \frac{2.4\times10^{-7}}{2.7\times10^{23}}
\approx 8.9\times10^{-31}\,\mathrm{W}.
\end{equation}
This is $\sim 10^{-37}$~W per Hz of bandwidth --- utterly undetectable.
\textbf{P15-based psi-communication is definitively infeasible at current
$\lambda$ bounds.} The Earth-Mars link requires modulating a natural
macroscopic mass.
\end{warning}

\subsection{Corrected Rate Table}

The W3i4 rate table is unchanged because $M_{\rm eff}$ does not enter:

\begin{table}[h]
\centering
\caption{Earth-Mars psi-channel capacity (W3i5 = W3i4, unchanged).
Source: planetary-mass transmitter $M_s = 8.3\times10^{23}$~kg.}
\begin{tabular}{lccc}
\hline\hline
Phase & $d$ (Mkm) & $\SNR_{\rm MRC}$ ($\eta=3$) & $C$ (kbits/s) \\
\hline
Opposition & $56$ & $14421$ & $14.0$ \\
Mean orbit & $225$ & $892$ & $10.2$ \\
Conjunction & $401$ & $281$ & $8.7$ \\
\hline\hline
\end{tabular}
\label{tab:earth-mars}
\end{table}

\confirmed{Earth-Mars rates unchanged by $M_{\rm eff}$ correction.
The rate $14.0$~kbits/s (opposition, MRC) requires $M_s \sim 10^{24}$~kg.}

%=============================================================================
\section{P8 Task 2: Does Inherent Secrecy Survive?}
\label{sec:secrecy}
%=============================================================================

\subsection{Secrecy Capacity Independence from $M_{\rm eff}$}

The psi-channel secrecy capacity (W3i4 Theorem~3) is:
\begin{equation}
C_s = B\log_2\!\left(\frac{1 + \SNR_B}{1 + \SNR_E}\right),
\label{eq:secrecy}
\end{equation}
where $\SNR_B$ is Bob's SNR (inside the corridor) and $\SNR_E$ is Eve's
SNR (outside the corridor). The key inequality $\SNR_B > \SNR_E$ follows
from:

\begin{enumerate}
\item The corridor geometry: Bob receives MRC-combined signal from all
  6~observables; Eve receives only the far-field scalar monopole.
\item The Hadamard bound: Eve's information is bounded by
  $I_{\rm Eve}/\text{bit} \leq 2\exp(-\pi d_E N / L)$, where $d_E$
  is Eve's distance from the corridor.
\end{enumerate}

Neither condition involves $M_{\rm eff}$. The secrecy capacity depends
on the \emph{ratio} $\SNR_B/\SNR_E$, which is a function of geometry
(corridor dimensions, Eve's position) only. Even if both $\SNR_B$ and
$\SNR_E$ scale identically with source mass (as they do), their
\emph{ratio} is mass-independent.

\begin{finding}[Secrecy Survives $M_{\rm eff}$ Correction]
\label{find:secrecy}
\confirmed{The inherent secrecy claim is \textbf{completely unaffected}
by the $M_{\rm eff}$ correction.} The secrecy capacity $C_s > 0$ for
all Eve positions outside the corridor, regardless of source mass,
mode mass, or absolute SNR level. Secrecy is a geometric property
of the psi-channel, not an amplitude property.

The only condition for $C_s > 0$ is that the channel must have
nonzero capacity to begin with: $\SNR_B > 0$. If there is any signal
at all, it is secret.
\end{finding}

%=============================================================================
\section{P8 Task 3: Corrected Channel Capacity}
\label{sec:capacity}
%=============================================================================

\subsection{SIMO Capacity --- Unchanged}

The W3i4-corrected capacity formula (SIMO MRC) is:
\begin{equation}
C_{\rm MRC} = B\log_2(1 + \SNR_{\rm eff}),
\quad \SNR_{\rm eff} = \sum_{j=1}^6 \frac{g_j^2 P_\xi}{\sigma_{n,j}^2}
\approx \eta \cdot \SNR_1, \quad \eta \approx 3.
\label{eq:C-MRC}
\end{equation}

The response coefficients $g_j$ scale as $G M_s / (c^2 d)$ --- they
use the \emph{transmitter mass}, not the mode mass. The noise variances
$\sigma_{n,j}^2$ are set by the receiver technology (atom interferometer
shot noise), not by the mode mass.

\begin{finding}[Channel Capacity Unaffected]
\label{find:capacity}
\confirmed{The SIMO MRC capacity $C_{\rm MRC}$ is unchanged by the
$M_{\rm eff}$ correction.} Laboratory capacity $\sim 98$~kbits/s
(if DFD is valid and a suitable source mass exists), Earth-Mars
opposition $\sim 14.0$~kbits/s (with planetary-mass transmitter).
\end{finding}

\subsection{Capacity Advantage Below $\sim 12$~dB SNR --- Unchanged}

The W3i4 comparison with 4D EM channels (Theorem~5) showed psi-channel
superiority only at $\SNR < s^* \approx 15$ (11.8~dB). This crossover
depends on the combining gain $\eta$ and the number of EM polarisation
modes (4), not on the mode mass. \confirmed{The SNR-regime comparison
is unchanged.}

%=============================================================================
\section{P8 Task 4: Honest Comparison with NASA DSN}
\label{sec:dsn}
%=============================================================================

\subsection{NASA Deep Space Network Performance}

The NASA Deep Space Network (DSN) communicates with Mars spacecraft
using:
\begin{itemize}
\item \textbf{Antennas}: 34~m (BWG) and 70~m (Mars station) dishes.
\item \textbf{Frequencies}: X-band (8.4~GHz uplink/downlink), Ka-band
  (32~GHz downlink).
\item \textbf{Mars opposition rates}: MRO Ka-band downlink achieves
  $\sim 5.2$~Mbits/s at closest approach ($\sim 56$~Mkm).
\item \textbf{Mars conjunction rates}: $\sim 0.5$--$10$~kbits/s
  (depending on Sun-Earth-probe angle; typically goes to safe mode
  within $\sim 2$--$5^\circ$ of the Sun).
\item \textbf{Conjunction blackout}: Complete blackout for $\sim 2$ weeks
  when the Sun-Earth-probe angle $< 2^\circ$.
\item \textbf{Power scaling}: $P_{\rm rec} \propto 1/d^2$ (free-space EM);
  data rate $\propto 1/d^2$ at fixed BER.
\end{itemize}

\subsection{DFD Psi-Channel vs DSN --- Head-to-Head}

\begin{table*}[t]
\centering
\caption{Honest comparison: DFD psi-channel vs NASA DSN for Earth-Mars
communication. All DFD numbers assume DFD physics is valid and a
planetary-mass transmitter ($M_s \sim 10^{24}$~kg) is available.}
\begin{tabular}{p{4.5cm}p{5.5cm}p{5.5cm}}
\hline\hline
\textbf{Parameter} & \textbf{DFD Psi-Channel} & \textbf{NASA DSN (X/Ka)} \\
\hline
Opposition data rate & 14.0~kbits/s (MRC) & $\sim 5.2$~Mbits/s (MRO Ka) \\
 & & \textbf{DSN wins: $370\times$ faster} \\
\hline
Conjunction data rate & 8.7~kbits/s (MRC) & $\sim 0$--$10$~kbits/s \\
 & (no blackout) & (blackout $\sim 2$ weeks) \\
 & & \textbf{DFD wins: no blackout} \\
\hline
Transmitter mass/power & $\sim 10^{24}$~kg (Venus-mass) & $\sim 100$~W TWTA, 35~kg HGA \\
 & & \textbf{DSN wins: $10^{22}\times$ lighter} \\
\hline
Power scaling with distance & $C \propto \log(1/d^2)$ & $C \propto 1/d^2$ \\
 & (logarithmic degradation) & (inverse-square degradation) \\
 & \textbf{DFD wins at large $d$} & \\
\hline
Secrecy & Inherent ($C_s > 0$, no key) & Requires encryption \\
 & \textbf{DFD wins} & \\
\hline
Conjunction transparency & $\alpha L_\odot = 4.6\times10^{-20}$ & Blocked by solar plasma \\
 & (effectively zero absorption) & ($> 60$~dB at $< 2^\circ$) \\
 & \textbf{DFD wins} & \\
\hline
Technology readiness & TRL~1 (theoretical) & TRL~9 (operational since 1964) \\
 & & \textbf{DSN wins} \\
\hline
Outer solar system & Voyager-equivalent at 150~AU: & Voyager 1 at 163~AU: \\
(Voyager-class distance) & $C \approx B\log_2(1 + \eta G^2 M_s^2 \xi_0^2$ & $\sim 160$~bits/s (S-band, 23~W) \\
 & $/(c^4 d^2 N_0 B))$ & \\
 & For $M_s = 10^{24}$~kg: & \\
 & $d = 2.2\times10^{13}$~m: & \\
 & $\SNR \approx 14421 \times (56/22000)^2 \approx 0.09$ & \\
 & $C \approx 130$~bits/s & \\
 & \textbf{Comparable, but requires Venus-mass TX} & \\
\hline\hline
\end{tabular}
\label{tab:dsn-comparison}
\end{table*}

\begin{finding}[Honest DSN Comparison]
\label{find:dsn}
\honest{The DFD psi-channel is \textbf{not competitive} with DSN on raw
data rate.} DSN delivers $370\times$ more data at Mars opposition with a
35~kg antenna, while DFD requires a Venus-mass transmitter for
$\sim 14$~kbits/s.

The DFD psi-channel has three genuine advantages over DSN:
\begin{enumerate}
\item \textbf{Zero conjunction blackout}: Continuous communication through
  solar conjunction, whereas DSN loses contact for $\sim 2$ weeks.
  For critical operations (Mars EDL during conjunction, emergency comms),
  this is operationally significant.
\item \textbf{Inherent secrecy}: No encryption overhead; interception
  is information-theoretically impossible outside the corridor.
\item \textbf{Logarithmic distance scaling}: At $d \gg d_{\rm Mars}$
  (outer solar system, interstellar), the DFD rate degrades as
  $\log(1/d^2)$ vs $1/d^2$ for EM. At Oort Cloud distances
  ($\sim 10^4$~AU), DFD retains $\sim 10$~bits/s while DSN approaches
  zero.
\end{enumerate}

The \emph{practical} value of DFD communications is therefore in
\textbf{niche deep-space scenarios} (conjunction-critical, ultra-deep
space, secure channels), not as a general replacement for EM. This is
a significant honest downgrade from the implication in earlier iterations
that DFD could ``replace'' deep-space EM communications.
\end{finding}

\begin{warning}[The Transmitter Mass Problem]
\label{warn:mass}
The entire P8 deep-space communications capability is contingent on
the ability to modulate a planetary-mass source at $\xi_0 = 10^{-3}$
and $f_c = 10^3$~Hz. No known or plausible technology can do this.
Even modulating an asteroid ($M \sim 10^{17}$~kg) at
$\xi_0 = 10^{-3}$ requires redistributing $10^{14}$~kg at kHz rates,
which is physically impossible with any foreseeable technology.

\killed{Deep-space psi-communication with useful data rates is
\textbf{technologically infeasible} at any plausible timescale.
The link budget requires transmitter masses that cannot be modulated.
P8 deep-space claims should be reclassified as \emph{theoretical}
(demonstrating that the physics permits it) rather than
\emph{technological} (demonstrating that it can be built).}
\end{warning}

%=============================================================================
\section{P8 Summary: What Survives W3i5}
\label{sec:p8-summary}
%=============================================================================

\begin{table}[h]
\centering
\caption{P8 W3i5 survival audit.}
\begin{tabular}{p{4.5cm}cc}
\hline\hline
Claim & Status & Ref. \\
\hline
Earth-Mars 14~kbits/s (MRC) & \confirmed{} (requires $10^{24}$~kg TX) & \S\ref{sec:earth-mars} \\
$M_{\rm eff}$ changes link budget & \corrected{No effect} & Finding~\ref{find:triage} \\
Secrecy $C_s > 0$ & \confirmed{Geometry-dependent, not mass-dependent} & \S\ref{sec:secrecy} \\
SIMO MRC capacity & \confirmed{Unchanged} & \S\ref{sec:capacity} \\
Psi beats DSN on rate & \killed{DSN $370\times$ faster} & Table~\ref{tab:dsn-comparison} \\
Conjunction transparency & \confirmed{} & \S\ref{sec:dsn} \\
P15 as psi-transmitter & \killed{$P_\psi \sim 10^{-31}$~W} & Warning~\ref{warn:p15-tx} \\
Practical deep-space comms & \killed{TX mass infeasible} & Warning~\ref{warn:mass} \\
\hline\hline
\end{tabular}
\label{tab:p8-summary}
\end{table}

%=============================================================================
%=============================================================================
\part*{Patent 9: Navigation and Positioning}
%=============================================================================
%=============================================================================

%=============================================================================
\section{P9 Task 5: $M_{\rm eff}$ Impact on Sub-mm Navigation}
\label{sec:p9-meff}
%=============================================================================

\subsection{P9 Positioning: Independent of $M_{\rm eff}$}

The P9 positioning system works by measuring the ambient $\psi$-field
gradient from natural masses (primarily Earth and nearby geological
structures). The position estimate is obtained by inverting the
measured field gradient:
\begin{equation}
\hat{\mathbf{r}} = \arg\min_{\mathbf{r}}
\left\|\nabla\psi_{\rm meas} - \nabla\psi_{\rm model}(\mathbf{r})\right\|^2,
\label{eq:p9-positioning}
\end{equation}
where $\psi_{\rm model}(\mathbf{r}) = G M_\oplus / (c^2 |\mathbf{r} - \mathbf{r}_\oplus|)$
plus geological perturbations.

The positioning precision is determined by:
\begin{equation}
\sigma_r = \frac{\sigma_\psi}{|\nabla^2\psi|} = \frac{\sigma_\psi\,c^2\,r^2}{2 G M_\oplus},
\label{eq:sigma-r}
\end{equation}
where $\sigma_\psi$ is the sensor noise floor (P6 quantum sensor).

\textbf{No dependence on $M_{\rm eff}$}: the field gradient $\nabla\psi$
is sourced by Earth's mass, not by a DFD cavity mode. The sensor noise
$\sigma_\psi$ is determined by atom interferometer performance. The DFD
mode mass enters only if one uses a \emph{cavity-based} sensor (P11 SQUID),
which is a P11 issue, not a P9 issue.

\begin{finding}[$M_{\rm eff}$ Does Not Affect P9 Positioning]
\label{find:p9-meff}
\confirmed{All P9 positioning claims are \textbf{completely independent}
of the $M_{\rm eff}$ correction.} The sub-mm positioning
($\sigma_r^{\rm HL} = 0.091$~mm underwater, W3i3) and the
sub-cm submarine navigation (W3i2) use natural gravitational
field gradients and P6 quantum sensors. Neither involves
the DFD cavity mode mass.
\end{finding}

\subsection{P9 with P11 SQUID Sensors --- Here $M_{\rm eff}$ Does Matter}

If a SQUID-based sensor (P11) were used instead of a P6 atom
interferometer, the SQUID couples to the psi-field via the cavity
mode, and its sensitivity is governed by $M_{\rm eff}$. The W3i4
correction makes SQUID-based psi-sensing $\sim 10^6\times$ less
sensitive than previously claimed.

\begin{correction}[P9+P11 SQUID Navigation --- Degraded]
\label{corr:p9-p11}
P9 with P11 SQUID sensors (if used for psi-field measurement):
the positioning precision degrades by $\sim 10^3\times$
(square root of the $10^6\times$ sensitivity degradation):
\begin{equation}
\sigma_r^{\rm P9+P11} \to \sigma_r^{\rm P9+P11,\,corrected}
\approx 10^3 \times \sigma_r^{\rm P9+P11,\,old}.
\end{equation}
This does not affect the baseline P9+P6 claims because P6
(atom interferometer) is the primary sensor pathway for P9.
\corrected{P9+P11 SQUID-based navigation is degraded by $\sim 10^3\times$
but this pathway was never the primary P9 design.}
\end{correction}

%=============================================================================
\section{P9 Task 6: Geological Detection Depth and Resolution}
\label{sec:geological}
%=============================================================================

\subsection{Independence from $M_{\rm eff}$}

The minimum detectable mass anomaly at depth $d$ (W3i4 Theorem~4):
\begin{equation}
\boxed{
M_{\rm anom}^{\min} = \frac{5\,\sigma_\psi^{\rm eff}\,c^2\,d}{G},
}
\label{eq:min-mass}
\end{equation}
involves:
\begin{itemize}
\item $\sigma_\psi^{\rm eff}$: sensor noise floor (P6 technology).
\item $c$, $G$: fundamental constants.
\item $d$: target depth.
\end{itemize}
No $M_{\rm eff}$ appears. The geological signal is $\delta\psi = G M_{\rm anom}/(c^2 d)$,
which is the Newtonian potential of the anomaly, entirely classical.

\begin{finding}[Geological Detection Unchanged]
\label{find:geological}
\confirmed{All W3i4 geological detection results survive the $M_{\rm eff}$
correction without modification:}
\begin{itemize}
\item 23~km depth, P6 HL sensors: $M_{\rm anom}^{\min} = 1.56\times10^{11}$~kg
  (840~m salt body).
\item 23~km depth, 8.7~nm P9+P6: $M_{\rm anom}^{\min} = 1.5\times10^{7}$~kg
  (15,000-tonne anomaly).
\item PSF lateral resolution: $d/\sqrt{2} = 16.3$~km (gradient PSF).
\item Vertical resolution: $\sigma_d \approx 0.6$~km at SNR~$= 5$.
\item Seismic noise discrimination: cross-spectral coherence $= 0$.
\end{itemize}
\end{finding}

%=============================================================================
\section{P9 Task 7: Does 8.7 nm Underwater Positioning Survive?}
\label{sec:87nm}
%=============================================================================

\subsection{Derivation Audit}

The 8.7~nm figure (W3i3, confirmed W3i4) arises from:
\begin{equation}
\sigma_r^{\rm P9+P6+array} = \frac{\sigma_r^{\rm P9+P6,\,HL}}
{\sqrt{N_{\rm sensors} \times N_{\rm meas}}}
= \frac{0.043\,\mathrm{mm}}{\sqrt{100 \times 10^3}}
\approx 8.7\,\mathrm{nm},
\label{eq:87nm}
\end{equation}
where:
\begin{itemize}
\item $\sigma_r^{\rm P9+P6,\,HL} = 0.043$~mm: Heisenberg-limited
  single-shot positioning from W3i3.
\item $N_{\rm sensors} = 100$: surface sensor array.
\item $N_{\rm meas} = T_{\rm int}/T_{\rm shot} = 1000$~s / $1$~s
  $= 10^3$: number of independent measurements.
\end{itemize}

\subsection{$M_{\rm eff}$ Dependence Check}

\begin{enumerate}
\item The $0.043$~mm Heisenberg-limited bound uses the P6 quantum
  Fisher information for measuring $\psi$ at a point, which depends
  on $N_{\rm photons}$ (atom number) and the field gradient $|\nabla\psi|$,
  not on $M_{\rm eff}$.
\item The $\sqrt{N_{\rm sensors} \times N_{\rm meas}}$ averaging is
  standard signal processing, independent of any DFD parameter.
\item The underwater environment provides $|\nabla\psi| \approx g/c^2
  \approx 1.1\times10^{-16}\,\mathrm{m}^{-1}$ from Earth's gravitational
  field, independent of $M_{\rm eff}$.
\end{enumerate}

\begin{finding}[8.7~nm Underwater Positioning Survives]
\label{find:87nm}
\confirmed{The 8.7~nm positioning claim is \textbf{completely unaffected}
by the $M_{\rm eff}$ correction.} The positioning precision depends on:
(a)~P6 quantum sensor noise floor, (b)~Earth's gravitational field
gradient, (c)~sensor count and integration time. None of these involve
the DFD cavity mode mass.

The strain sensitivity $\varepsilon = 8.7\times10^{-9}$ (1~m baseline)
and the seismograph capability ($M \geq 6.5$ earthquakes) are likewise
unchanged.
\end{finding}

\subsection{Caveats on the 8.7~nm Claim}

While the $M_{\rm eff}$ correction does not affect the number, several
assumptions deserve honest scrutiny:

\begin{enumerate}
\item \textbf{Heisenberg-limited sensors}: The $0.043$~mm single-shot
  bound assumes NOON-state entangled atoms achieving the Heisenberg limit.
  Current atom interferometers operate at $\sim 10\times$~SQL, not HL.
  The SQL-limited 8.7~nm would become:
  \begin{equation}
  \sigma_r^{\rm SQL} = 8.7\,\mathrm{nm} \times \sqrt{N_{\rm atom}}
  \approx 8.7\,\mathrm{nm} \times 10^3 = 8.7\,\mu\mathrm{m}.
  \end{equation}
  Still excellent ($< 10\,\mu$m underwater positioning) but not sub-nm.

\item \textbf{Coherent averaging}: The $\sqrt{10^5}$ averaging factor
  assumes all 100 sensors maintain phase coherence over 1000~s. Phase
  coherence requires P5 timing at $\sigma_t < 10^{-12}$~s (W3i4 P9
  cross-reference). If timing coherence is limited to $T_{\rm coh} < 1000$~s,
  the averaging gain is reduced.

\item \textbf{Systematic floors}: Tidal loading ($\pm 4$~mm, 12.4~hr period),
  ocean loading, atmospheric pressure, and nearby mass redistribution
  (vehicles, tides) create systematic psi-field variations that must be
  modelled and subtracted. The 8.7~nm figure assumes perfect subtraction.
\end{enumerate}

%=============================================================================
\section{P9 Task 8: Impact of $n_\psi = \sqrt{2}$ on MOND-Zone Positioning}
\label{sec:npsi}
%=============================================================================

\subsection{The $n_\psi = \sqrt{2}$ Result (Recalled)}

W3i4 P12 proved that the psi-field refractive index approaches
$n_\psi \to \sqrt{2}$ universally in the deep MOND limit
($|\nabla\psi| \ll a_0$), valid for any interpolating function
satisfying $W(y) \propto y^{3/4}$ for $y \ll 1$. The previous
assumption $n_\psi \to \infty$ is superseded.

\subsection{Where Does $n_\psi$ Enter P9?}

The P9 positioning system measures the \emph{static} psi-field
$\psi(\mathbf{r}) = GM/(c^2 r)$, which is a DC (zero-frequency)
quantity. The refractive index $n_\psi$ governs the \emph{propagation
speed} of psi-field \emph{perturbations} (waves), not the static field
value.

For \emph{ranging-based} positioning (measuring propagation time of
psi-pulses between beacons), the effective speed of propagation is:
\begin{equation}
c_s = \frac{c}{n_\psi(|\nabla\psi|)}.
\label{eq:cs}
\end{equation}

In the Newtonian regime ($|\nabla\psi| \gg a_0$, i.e., near Earth's
surface where $g \approx 9.8$~m/s$^2 \gg a_0 \approx 1.2\times10^{-10}$~m/s$^2$):
\begin{equation}
n_\psi \approx 1 + \order{a_0/g} \approx 1 + 10^{-11}.
\label{eq:npsi-surface}
\end{equation}
The correction to surface-level positioning is $\sim 10^{-11}$,
corresponding to $\sim 0.06$~pm per metre of baseline --- utterly
negligible compared to the 8.7~nm precision.

\subsection{MOND-Zone Positioning: Altitude $\gtrsim 3200$~km}

The MOND acceleration threshold is crossed at altitude $h$ where
$g(R_\oplus + h) = a_0$:
\begin{equation}
h_{\rm MOND} = R_\oplus\!\left(\sqrt{\frac{g_0}{a_0}} - 1\right)
\approx 6370\,\mathrm{km}\!\left(\sqrt{\frac{9.8}{1.2\times10^{-10}}} - 1\right)
\approx 6370\!\left(2.86\times10^5 - 1\right)\,\mathrm{km}.
\label{eq:h-mond}
\end{equation}
This gives $h_{\rm MOND} \approx 1.82\times10^9$~km, which is
$\sim 12$~AU --- well outside the solar system for Earth-orbiting
applications. The ``3200~km'' altitude for MOND-zone effects quoted
in earlier iterations refers to the altitude where the \emph{galactic}
gradient dominates over the \emph{solar} gradient, not where $g < a_0$.

\begin{correction}[MOND-Zone Altitude for P9]
\label{corr:mond-zone}
\corrected{The MOND zone ($|\nabla\psi| < a_0$) is reached not at
3200~km altitude but at $\sim 12$~AU from the Sun (where the solar
gravitational acceleration drops below $a_0$).} At 3200~km altitude,
$g \approx 4.3$~m/s$^2$, which is $3.6\times10^{10}$ times above $a_0$.
The $n_\psi = \sqrt{2}$ correction applies only in truly deep MOND
environments: interstellar space, galaxy outskirts, or voids.

The ``3200~km'' figure from W3i4 P12 refers to the psi-ionosphere
(where the \emph{galactic} $\psi$-gradient creates total internal
reflection for psi-waves), not the MOND acceleration threshold.
These are distinct physical scales.
\end{correction}

\subsection{Impact on P9 Positioning: Quantitative}

\begin{finding}[$n_\psi = \sqrt{2}$ Impact on P9]
\label{find:npsi-p9}
For the three P9 operating regimes:

\textbf{1. Surface/underwater navigation} (primary use case):
$n_\psi - 1 \lesssim 10^{-11}$. Effect on positioning: $< 0.1$~pm.
\confirmed{Negligible. No correction required.}

\textbf{2. LEO/MEO satellite navigation} ($h = 200$--$36000$~km):
$g \approx 0.2$--$9$~m/s$^2$, still $\gg a_0$.
$n_\psi - 1 \lesssim 10^{-10}$. Effect on ranging: $< 1$~pm.
\confirmed{Negligible.}

\textbf{3. Deep-space navigation} ($d > 10$~AU from Sun):
Here $g_\odot < 10^{-4}$~m/s$^2$ and the galactic field
$g_{\rm gal} \sim 2\times10^{-10}$~m/s$^2 \approx a_0$.
$n_\psi \to \sqrt{2}$, giving $c_s \to c/\sqrt{2}$.
Ranging-based positioning must use $c_s = c/\sqrt{2}$, not $c$.
This is a $\sqrt{2} - 1 \approx 41\%$ systematic error if uncorrected.

For P9 deep-space positioning at $d = 100$~AU:
\begin{equation}
\Delta r_{\rm syst} = r \times (n_\psi - 1)/n_\psi
\approx r \times 0.293 \approx 4.4\times10^{12}\,\mathrm{m} \times 0.293
\approx 1.3\times10^{12}\,\mathrm{m} \approx 8.7\,\mathrm{AU}.
\label{eq:deep-space-error}
\end{equation}
\caution{If psi-pulse ranging is used for deep-space positioning without
correcting for $n_\psi = \sqrt{2}$, the systematic position error is
$\sim 29\%$ of the distance --- catastrophically large. However, this is
a \emph{known} correction: once $n_\psi(\mathbf{r})$ is modelled, the
systematic is removed. The residual error depends on the accuracy of
the $n_\psi$ model, estimated at $\sim 1\%$ of the correction
$\approx 0.3\%$ of distance.}
\end{finding}

\subsection{$n_\psi = \sqrt{2}$ Effect on Geological Sensing}

The geological PSF (W3i4 Sec.~III) uses the \emph{static} field
$\delta\psi = G M_{\rm anom}/(c^2 d)$, not psi-wave propagation.
The refractive index governs wave propagation speed and does not
affect static field measurements.

\confirmed{Geological detection PSF, lateral resolution, minimum
detectable anomaly, and all inversion results are \textbf{completely
unaffected} by $n_\psi = \sqrt{2}$.}

%=============================================================================
\section{P9+P6 Strain Sensitivity: Corrected Value}
\label{sec:strain}
%=============================================================================

\subsection{The W3i4 Claim: $5.8\times10^{-17}$~m/m}

The iteration tracker reports ``P9+P6 combined: strain sensitivity
$5.8\times10^{-17}$~m/m.'' This value appears to have been derived
from a different context (psi-field frequency shift, not position-based
strain). The W3i4 P9 text (Eq.~30) derives strain from 8.7~nm
positioning:
\begin{equation}
\varepsilon = \frac{\sigma_r}{L_{\rm base}} = \frac{8.7\times10^{-9}}{L_{\rm base}}.
\end{equation}

For $L_{\rm base} = 1$~m: $\varepsilon = 8.7\times10^{-9}$ (consistent
with W3i4 P9 Finding~5).

For $L_{\rm base} = 150$~km (continental scale):
$\varepsilon = 8.7\times10^{-9}/1.5\times10^5 = 5.8\times10^{-14}$~m/m.

The $5.8\times10^{-17}$~m/m figure in the tracker would require
$L_{\rm base} = 150,000$~km $\approx 0.001$~AU, which is not a
realistic terrestrial baseline. This value likely arises from a
different calculation (perhaps the psi-field frequency sensitivity,
not position-based strain).

\begin{correction}[P9+P6 Strain Sensitivity]
\label{corr:strain}
The \textbf{position-based} strain sensitivity from 8.7~nm differential
positioning is:
\begin{center}
\begin{tabular}{cc}
\toprule
Baseline & Strain sensitivity \\
\midrule
1~m & $8.7\times10^{-9}$ \\
10~m & $8.7\times10^{-10}$ \\
100~m & $8.7\times10^{-11}$ \\
1~km & $8.7\times10^{-12}$ \\
\bottomrule
\end{tabular}
\end{center}
\corrected{The ``$5.8\times10^{-17}$~m/m'' figure in the iteration
tracker requires clarification of the baseline assumed. For any
reasonable terrestrial baseline ($L \leq 10$~km), the strain
sensitivity is $\varepsilon \geq 8.7\times10^{-13}$.}
\end{correction}

%=============================================================================
\section{T1 Falsification and T5 Dual-Role: Impact on P8/P9}
\label{sec:tensions}
%=============================================================================

\subsection{T1: $\dot{G}/G$ Falsification}

W3i4 P5/P14 declared T1 (DFD predicts $\dot{G}/G = 2.6\times10^{-11}$~yr$^{-1}$,
violating LLR bound $< 7.5\times10^{-14}$ by $347\times$) as a genuine
falsification of the \emph{cosmological} sector of DFD.

\textbf{Impact on P8}: The P8 link budget uses $G$ as a constant.
If $G$ varies at $2.6\times10^{-11}$~yr$^{-1}$, the effect on the
link budget over a Mars synodic period ($\sim 2.14$~yr) is:
$\Delta G/G \approx 5.6\times10^{-11}$ --- negligible for all P8 purposes.
\confirmed{T1 does not affect P8 communications.}

\textbf{Impact on P9}: The P9 positioning model uses $\psi = GM/(c^2 r)$.
A time-varying $G$ causes a slow drift in $\psi$:
$\dot{\psi}/\psi = \dot{G}/G = 2.6\times10^{-11}$~yr$^{-1}$.
Over 1000~s integration: $\Delta\psi/\psi \approx 8.2\times10^{-19}$
--- negligible compared to the sensor noise floor ($10^{-14}$~per shot).
\confirmed{T1 does not affect P9 positioning at any practical timescale.}

However: T1 is a theoretical problem for DFD as a whole. If the
cosmological sector is falsified, the \emph{physical validity of the
psi-field itself} is in question. P8 and P9 rely on DFD being physically
correct. T1 threatens the entire edifice, not just the cosmological
predictions.

\subsection{T5: Psi-Field Dual Role}

W3i4 identified T5: the psi-field simultaneously plays the role of
(a)~the Newtonian gravitational potential and (b)~an EM-sourced scalar
field. These two roles are structurally inconsistent in the DFD field
equations.

\textbf{Impact on P9}: P9 positioning relies on the psi-field being
the Newtonian potential ($\psi = GM/(c^2 r)$). If T5 is resolved by
separating $\psi$ into two distinct fields ($\psi_{\rm grav}$ and
$\psi_{\rm EM}$), then P9 positioning would use $\psi_{\rm grav}$ and
would be entirely conventional gravitational navigation. The DFD-specific
features (nonreciprocal timing, MOND-zone corrections) would involve
$\psi_{\rm EM}$.

\textbf{Impact on P8}: If $\psi$ is separated, the psi-channel (P8) would
carry modulations of $\psi_{\rm EM}$, not $\psi_{\rm grav}$. The
coupling strength would change depending on which field carries the
information. This is an open theoretical question that cannot be
resolved until T5 is addressed at the field-equation level.

\caution{T5 resolution could fundamentally change the physical content
of both P8 and P9. Until T5 is resolved, all P8/P9 predictions carry
a theoretical uncertainty associated with the ambiguity in $\psi$.}

%=============================================================================
\section{Top 5 New Findings (W3i5, P8+P9 Combined)}
\label{sec:top5}
%=============================================================================

\begin{center}
\begin{tabular}{clcc}
\toprule
\# & Finding & Category & Impact \\
\midrule
1 & P8: DSN is $370\times$ faster than DFD at Mars;
  DFD advantage is conjunction/secrecy/deep-space only
  & Honest audit & \textbf{Critical} \\
2 & P8: P15 as psi-transmitter gives $P_\psi \sim 10^{-31}$~W;
  deep-space psi-comms is technologically infeasible
  & \textbf{KILLED} & \textbf{Critical} \\
3 & P8+P9: $M_{\rm eff}$ correction has \textbf{zero effect}
  on free-space propagation and passive sensing
  & Correction scope & High \\
4 & P9: MOND zone at 12~AU, not 3200~km; $n_\psi = \sqrt{2}$
  is irrelevant for terrestrial/LEO positioning
  & Correction & High \\
5 & P9: 8.7~nm underwater positioning and all geological
  detection results survive without modification
  & Survival audit & High \\
\bottomrule
\end{tabular}
\end{center}

%=============================================================================
\section{Open Questions for W3i6}
\label{sec:open}
%=============================================================================

\begin{enumerate}
\item \textbf{P8: Near-field psi-communication}: The deep-space link
  is infeasible, but what about short-range ($< 1$~km) communication
  using laboratory-scale masses ($M_s = 10^3$--$10^6$~kg)? Derive the
  maximum range for 1~bit/s at $M_s = 10^6$~kg. This is the only
  physically realisable P8 scenario.

\item \textbf{P8: Passive eavesdropping detection}: Can the corridor
  geometry detect an eavesdropper via their mass perturbation on the
  psi-field? What minimum Eve mass is detectable?

\item \textbf{P9: Time-lapse geological monitoring}: Continuous
  psi-sensor networks for aquifer depletion or magma chamber monitoring.
  What temporal resolution is achievable for $\Delta\rho \sim 10$~kg/m$^3$
  changes at 5~km depth?

\item \textbf{P9: Impact of T5 resolution}: If T5 is resolved by
  field splitting ($\psi_{\rm grav} \neq \psi_{\rm EM}$), which field
  does P9 positioning use? Does the MOND correction survive?

\item \textbf{P8+P9 combined}: Can P9's geological sensing capability
  be used to \emph{locate} natural mass concentrations suitable as
  P8 psi-transmitters? Derive the minimum geological anomaly that
  could serve as a modulated psi-source (e.g., by periodic water
  injection/extraction).
\end{enumerate}

%=============================================================================
\section{Dimensional Consistency Checks}
\label{sec:dimensions}
%=============================================================================

\paragraph{Link budget (Eq.~\ref{eq:link-budget}):}
$[G^2 M_s^2 \xi_0^2/(c^4 d^2 N_0 B)]
= [\mathrm{m}^6\,\mathrm{kg}^{-2}\,\mathrm{s}^{-4}
\cdot \mathrm{kg}^2 \cdot 1]/[\mathrm{m}^4\,\mathrm{s}^{-4}
\cdot \mathrm{m}^2 \cdot \mathrm{m}^2\,\mathrm{s}^{-2}\,\mathrm{Hz}^{-1}
\cdot \mathrm{Hz}] = 1$. \checkmark (dimensionless SNR)

\paragraph{Minimum detectable mass (Eq.~\ref{eq:min-mass}):}
$[5\sigma_\psi c^2 d / G]
= [1 \cdot \mathrm{m}^2\,\mathrm{s}^{-2} \cdot \mathrm{m}
/ (\mathrm{m}^3\,\mathrm{kg}^{-1}\,\mathrm{s}^{-2})]
= \mathrm{kg}$. \checkmark

\paragraph{Positioning precision (Eq.~\ref{eq:sigma-r}):}
$[\sigma_\psi c^2 r^2 / (G M)]
= [1 \cdot \mathrm{m}^2\,\mathrm{s}^{-2} \cdot \mathrm{m}^2
/ (\mathrm{m}^3\,\mathrm{kg}^{-1}\,\mathrm{s}^{-2} \cdot \mathrm{kg})]
= \mathrm{m}$. \checkmark

%=============================================================================
\begin{thebibliography}{99}

\bibitem{W3i4P8}
DFD Research Programme,
``W3i4 P8: 6D Shannon Capacity Formal Proof, Earth-Mars Link Budget,
QKD Rate Analysis,'' March 2026.

\bibitem{W3i4P9}
DFD Research Programme,
``W3i4 P9: Geological Detection Resolution Function, PSF Derivation,
8.7~nm Positioning and Strain Sensing,'' March 2026.

\bibitem{W3i4P11}
DFD Research Programme,
``W3i4 P11: Mode Mass Resolution, SQUID Sensitivity Correction,''
March 2026.

\bibitem{W3i4P12}
DFD Research Programme,
``W3i4 P12: $n_\psi = \sqrt{2}$ Universal Proof, GRIN Lens,''
March 2026.

\bibitem{DSN810-005}
NASA/JPL, ``DSN Telecommunications Link Design Handbook,''
Document 810-005, Rev.~E, 2022.

\bibitem{MRO-telecom}
A.~Makovsky, A.~Barbieri, R.~Tung,
``Mars Reconnaissance Orbiter Telecommunications,''
DESCANSO Design and Performance Summary Series, Article 12, 2009.

\end{thebibliography}

\end{document}
